% SOURCE_AUTHORITY: sga5_fr_workpass.tex, lines 6878--6924 (LF-normalized unit).
% SOURCE_SHA256: 791F4EFFC5E02832D5D77ED03518C8156D6F07E4C8238B03545DB93D883FBB28
\subsubsection*{2.4. A categoria $P_{AR}=\underline{\Hom}_{AR}(\NN^\circ,\cC)$}
Retomando as notações de (2.1), seja $AR$ o conjunto das flechas canônicas (cf. 2.1)
\[
v_{rX}:X[r]\longrightarrow X,
\]
onde $X$ percorre o conjunto de objetos de $P$ e $r$ o conjunto de inteiros $\ge0$.

\begin{statement}{Proposição 2.4.1}
O subconjunto $AR$ de $\Fl(P)$ permite simultaneamente um cálculo de frações pela direita e um cálculo de frações pela esquerda.
\end{statement}

\begin{prooftext}[Demonstração]
As implicações (a) $\Rightarrow$ (a)' e (b) $\Rightarrow$ (b)' de (2.3.1) e (2.3.2) são claras.

Demonstremos (c). Seja, pois, $(D_1)$ o diagrama comutativo que aparece a seguir, no qual $v_{rY}$ designa a flecha canônica de $Y[r]$ em $Y$:
\[
\begin{tikzcd}[column sep=large,row sep=large]
& Y[r] \arrow[d,"{v_{rY}}"]\\
X \arrow[r] & Y .
\end{tikzcd}
\]
O diagrama $(D_2)$ que aparece a seguir é comutativo e $v_{rX}$, morfismo canônico de $X[r]$ em $X$, pertence a $AR$, de onde segue a afirmação:
\[
\begin{tikzcd}[column sep=large,row sep=large]
X[r] \arrow[r,"{u[r]}"] \arrow[d,"{v_{rX}}"'] & Y[r] \arrow[d,"{v_{rY}}"]\\
X \arrow[r,"u"'] & Y .
\end{tikzcd}
\]
Demonstremos (d). Seja um diagrama comutativo:
\[
\begin{tikzcd}[column sep=large]
X \arrow[r,shift left,"f"] \arrow[r,shift right,"g"'] & Z[r] \arrow[r,"{v_{rZ}}"] & Z .
\end{tikzcd}
\]
Aplicando a translação por $r$, deduz-se um diagrama comutativo
\[
\begin{tikzcd}[column sep=large]
X[r] \arrow[r,shift left,"{f[r]}"] \arrow[r,shift right,"{g[r]}"'] & Z[2r] \arrow[r,"{v_{rZ}[r]}"] & Z[r].
\end{tikzcd}
\]
É claro que o diagrama que aparece a seguir é comutativo, de modo que $f\circ v_{rX}=g\circ v_{rX}$:
\[
\begin{tikzcd}[column sep=large,row sep=large]
X[r] \arrow[r,shift left,"{f[r]}"] \arrow[r,shift right,"{g[r]}"'] \arrow[d,"{v_{rX}}"'] & Z[2r] \arrow[d,"{v_{rZ}[r]}"]\\
X \arrow[r,shift left,"f"] \arrow[r,shift right,"g"'] & Z[r] .
\end{tikzcd}
\]
